\documentclass[12pt, a4paper]{article}

% ============================================
% КОДИРОВКА И ЯЗЫК
% ============================================
\usepackage[utf8]{inputenc}
\usepackage[T2A]{fontenc}
\usepackage[russian, english]{babel}

% ============================================
% ГЕОМЕТРИЯ СТРАНИЦЫ
% ============================================
\usepackage[
  left=3cm,
  right=1.5cm,
  top=2cm,
  bottom=2cm
]{geometry}

% ============================================
% БИБЛИОГРАФИЯ
% ============================================
\usepackage[
  backend=biber,
  style=gost-numeric,      % Можно изменить на: authoryear, numeric, alphabetic, apa
  sorting=none,            % none - по порядку цитирования, nyt - имя-год-название
  maxbibnames=99,          % Показывать всех авторов в библиографии
  maxcitenames=2,          % Максимум 2 автора в цитировании
  giveninits=true,         % Инициалы вместо полных имен
  doi=true,                % Показывать DOI
  isbn=true,               % Показывать ISBN
  url=true,                % Показывать URL
  eprint=false             % Не показывать eprint
]{biblatex}

% Подключение файла библиографии
\addbibresource{references.bib}

% Настройка заголовка библиографии
\DefineBibliographyStrings{russian}{
  bibliography = {Список использованных источников},
  references = {Ссылки},
}

% ============================================
% ДОПОЛНИТЕЛЬНЫЕ ПАКЕТЫ
% ============================================

% Графика
\usepackage{graphicx}
\graphicspath{{images/}}

% Цвета
\usepackage{xcolor}
\definecolor{linkcolor}{RGB}{0, 0, 180}
\definecolor{urlcolor}{RGB}{0, 128, 0}
\definecolor{citecolor}{RGB}{180, 0, 0}

% Гиперссылки
\usepackage{hyperref}
\hypersetup{
  colorlinks=true,
  linkcolor=linkcolor,
  urlcolor=urlcolor,
  citecolor=citecolor,
  bookmarksnumbered=true,
  pdftitle={Отчет по работе с библиографией в LaTeX},
  pdfauthor={Ваше Имя},
  pdfsubject={Лабораторная работа},
  pdfkeywords={LaTeX, BibTeX, библиография}
}

% Правильные кавычки
\usepackage{csquotes}

% Листинги кода
\usepackage{listings}
\lstset{
  basicstyle=\ttfamily\small,
  breaklines=true,
  frame=single,
  numbers=left,
  numberstyle=\tiny,
  backgroundcolor=\color{gray!10},
  commentstyle=\color{green!50!black},
  keywordstyle=\color{blue},
  stringstyle=\color{red!50!black},
  showstringspaces=false,
  tabsize=2
}

% Таблицы
\usepackage{booktabs}
\usepackage{longtable}
\usepackage{array}

% Математика
\usepackage{amsmath, amssymb}

% Подписи к рисункам и таблицам
\usepackage{caption}
\captionsetup{
  labelsep=period,
  justification=centering,
  font=small,
  labelfont=bf
}

% Колонтитулы
\usepackage{fancyhdr}
\pagestyle{fancy}
\fancyhf{}
\fancyhead[L]{\small Лабораторная работа}
\fancyhead[R]{\small Работа с библиографией}
\fancyfoot[C]{\thepage}
\renewcommand{\headrulewidth}{0.4pt}
\renewcommand{\footrulewidth}{0pt}

% Интервалы
\usepackage{setspace}
\onehalfspacing  % Полуторный интервал

% Отступ первой строки
\usepackage{indentfirst}
\setlength{\parindent}{1.25cm}

% ============================================
% ИНФОРМАЦИЯ О ДОКУМЕНТЕ
% ============================================
\title{
  \textbf{Лабораторная работа №6} \\
  \vspace{0.5cm}
  \Large Работа с библиографией в LaTeX \\
  \large Использование BibTeX и biblatex
}

\author{
  Выполнил: Студент Группы НПМмд02-24 Викторов Егор \\
}

\date{\today}

% ============================================
% НАЧАЛО ДОКУМЕНТА
% ============================================
\begin{document}

% Титульная страница
\begin{titlepage}
  \centering
  
  \vspace*{2cm}
  
  {\LARGE \textbf{МИНИСТЕРСТВО НАУКИ И ВЫСШЕГО ОБРАЗОВАНИЯ} \\
  \textbf{РОССИЙСКОЙ ФЕДЕРАЦИИ}}
  
  \vspace{0.5cm}
  
  {\large Федеральное государственное бюджетное образовательное учреждение \\
  высшего образования}
  
  \vspace{0.3cm}
  
  {\Large \textbf{«РУДН»}}
  
  \vspace{0.5cm}
  
  {\large Факультет физико-математических и естественных наук \\
  Кафедра Теории Вероятностей и Кибербезопасности}
  
  \vspace{3cm}
  
  {\LARGE \textbf{ЛАБОРАТОРНАЯ РАБОТА №6}}
  
  \vspace{0.5cm}
  
  {\Large по дисциплине «Компьютерный практикум по научному письму»}
  
  \vspace{0.5cm}
  
  {\large на тему:}
  
  \vspace{0.3cm}
  
  {\Large \textbf{«Работа с библиографией в LaTeX. \\
  Использование BibTeX и biblatex»}}
  
  \vfill
  
  \begin{flushright}
    \begin{tabular}{ll}
      Выполнил: & студент группы НПМмд02-24 \\
                & Викторов Е.И. \\
    \end{tabular}
  \end{flushright}
  
  \vspace{2cm}
  
  {\large Москва \\ \the\year}
  
\end{titlepage}

% Содержание
\tableofcontents
\newpage

% ============================================
% ВВЕДЕНИЕ
% ============================================
\section{Введение}

Управление библиографией является важной частью подготовки научных и технических документов. Система LaTeX предоставляет мощные инструменты для автоматизации этого процесса через пакеты BibTeX и biblatex \parencite{latex-project}.

\subsection{Цель работы}

Целью данной лабораторной работы является изучение методов работы с библиографическими ссылками в системе LaTeX, освоение инструментов BibTeX и biblatex для автоматизации создания списка литературы.

\subsection{Задачи работы}

В рамках лабораторной работы необходимо решить следующие задачи:

\begin{enumerate}
  \item Изучить структуру файлов библиографии формата \texttt{.bib}
  \item Освоить различные типы библиографических записей
  \item Научиться использовать команды цитирования в тексте документа
  \item Изучить различные стили оформления библиографии
  \item Настроить автоматическую компиляцию документа с библиографией
  \item Создать полноценный документ с корректно оформленными ссылками
\end{enumerate}

\subsection{Актуальность темы}

Как отмечает \textcite{lamport1994}, правильное оформление библиографических ссылок является обязательным требованием для научных публикаций. Автоматизация этого процесса позволяет исследователям сосредоточиться на содержании работы, а не на технических деталях оформления.

% ============================================
% ТЕОРЕТИЧЕСКАЯ ЧАСТЬ
% ============================================
\section{Теоретическая часть}

\subsection{История развития систем управления библиографией}

Система TeX была создана Дональдом Кнутом в конце 1970-х годов \parencite{knuth1984}. Одним из ключевых алгоритмов, лежащих в основе TeX, является алгоритм разбиения абзацев на строки, описанный в работе \textcite{knuth1977}.

Позже Лесли Лампорт разработал надстройку над TeX, получившую название LaTeX, которая значительно упростила создание структурированных документов \parencite{lamport1994}.

\subsection{Инструменты для работы с библиографией}

\subsubsection{BibTeX}

BibTeX — это программа и формат файлов для управления библиографическими ссылками, разработанная Ореном Паташником \parencite{patashnik1988}. Основные характеристики:

\begin{itemize}
  \item Простой текстовый формат хранения библиографических данных
  \item Поддержка различных типов источников (книги, статьи, веб-ресурсы)
  \item Автоматическая сортировка и форматирование списка литературы
  \item Широкая поддержка стилей оформления
\end{itemize}

Подробное описание создания стилей для BibTeX приведено в техническом отчете \textcite{bibtex-guide}.

\subsubsection{Biblatex и Biber}

Biblatex — это современная альтернатива BibTeX, предоставляющая расширенные возможности \parencite{biblatex-manual}:

\begin{itemize}
  \item Полная поддержка Unicode
  \item Гибкая настройка стилей цитирования
  \item Расширенные возможности сортировки
  \item Поддержка локализации на множество языков
  \item Использование Biber в качестве backend-процессора \parencite{biber-manual}
\end{itemize}

\subsection{Структура файла библиографии}

Файл библиографии с расширением \texttt{.bib} содержит записи различных типов. Каждая запись имеет следующую структуру:

\begin{lstlisting}[language=TeX, caption={Общая структура BibTeX записи}]
@ТИП_ИСТОЧНИКА{ключ_цитирования,
  поле1 = {значение1},
  поле2 = {значение2},
  поле3 = {значение3}
}
\end{lstlisting}

\subsubsection{Типы библиографических записей}

В таблице \ref{tab:bibtex-types} представлены основные типы записей BibTeX.

\begin{table}[htbp]
  \centering
  \caption{Основные типы BibTeX записей}
  \label{tab:bibtex-types}
  \begin{tabular}{@{}lp{8cm}@{}}
    \toprule
    \textbf{Тип} & \textbf{Описание} \\
    \midrule
    \texttt{@article} & Статья в журнале или периодическом издании \\
    \texttt{@book} & Книга с указанием издателя \\
    \texttt{@inproceedings} & Статья в материалах конференции \\
    \texttt{@online} & Интернет-ресурс \\
    \texttt{@techreport} & Технический отчет организации \\
    \texttt{@manual} & Техническая документация \\
    \texttt{@misc} & Источник, не подходящий под другие категории \\
    \texttt{@phdthesis} & Диссертация \\
    \texttt{@mastersthesis} & Магистерская диссертация \\
    \bottomrule
  \end{tabular}
\end{table}

\subsection{Команды цитирования}

Пакет biblatex предоставляет множество команд для цитирования источников. Основные команды представлены в таблице \ref{tab:cite-commands}.

\begin{table}[htbp]
  \centering
  \caption{Команды цитирования в biblatex}
  \label{tab:cite-commands}
  \begin{tabular}{@{}llp{5cm}@{}}
    \toprule
    \textbf{Команда} & \textbf{Пример результата} & \textbf{Использование} \\
    \midrule
    \verb|\cite{key}| & [1] & Базовое цитирование \\
    \verb|\parencite{key}| & [1] & Цитирование в скобках \\
    \verb|\textcite{key}| & Knuth [1] & Текстовое цитирование \\
    \verb|\citeauthor{key}| & Knuth & Только автор \\
    \verb|\citeyear{key}| & 1984 & Только год \\
    \verb|\citetitle{key}| & The TeXbook & Только название \\
    \verb|\footcite{key}| & \textsuperscript{1} & Цитирование в сноске \\
    \bottomrule
  \end{tabular}
\end{table}

% ============================================
% ПРАКТИЧЕСКАЯ ЧАСТЬ
% ============================================
\section{Практическая часть}

\subsection{Создание файла библиографии}

Для выполнения лабораторной работы был создан файл \texttt{references.bib}, содержащий различные типы библиографических записей. Файл включает:

\begin{itemize}
  \item 3 записи типа \texttt{@book} — классические книги по LaTeX
  \item 2 записи типа \texttt{@article} — научные статьи
  \item 5 записей типа \texttt{@online} — интернет-ресурсы
  \item 1 запись типа \texttt{@inproceedings} — материалы конференции
  \item 1 запись типа \texttt{@techreport} — технический отчет
\end{itemize}

\subsection{Настройка документа}

В преамбуле документа был подключен пакет biblatex с следующими параметрами:

\begin{lstlisting}[language=TeX, caption={Настройка biblatex}]
\usepackage[
  backend=biber,
  style=gost-numeric,
  sorting=none,
  maxbibnames=99,
  maxcitenames=2,
  giveninits=true,
  doi=true,
  isbn=true,
  url=true
]{biblatex}
\addbibresource{references.bib}
\end{lstlisting}

\subsection{Примеры цитирования}

\subsubsection{Цитирование книг}

Классическая книга по TeX \parencite{knuth1984} является фундаментальным трудом в области компьютерной типографики. Для изучения LaTeX рекомендуется книга \textcite{lamport1994}, а для углубленного изучения — работа \textcite{mittelbach2004}.

\subsubsection{Цитирование статей}

Историческая статья \textcite{dijkstra1968} оказала значительное влияние на развитие структурного программирования. Алгоритмические основы разбиения текста на строки описаны в работе \parencite{knuth1977}.

\subsubsection{Цитирование онлайн-ресурсов}

Официальная документация LaTeX доступна на сайте проекта \parencite{latex-project}. Для начинающих пользователей рекомендуется учебное пособие \parencite{overleaf-docs}. Пакеты для LaTeX можно найти в репозитории CTAN \parencite{ctan}.

Русскоязычные пользователи могут обратиться к ресурсам \parencite{latex-ru} и книге \textcite{lvovsky2003}.

\subsubsection{Множественное цитирование}

Основы работы с LaTeX описаны во многих источниках \parencite{knuth1984, lamport1994, mittelbach2004}. Для решения практических задач полезны онлайн-ресурсы \parencite{overleaf-docs, latex-project, tex-stackexchange}.

\subsection{Процесс компиляции}

Для корректной работы библиографии документ необходимо компилировать в следующей последовательности:

\begin{enumerate}
  \item \texttt{pdflatex report.tex} — первая компиляция, создание \texttt{.aux} файла
  \item \texttt{biber report} — обработка библиографии
  \item \texttt{pdflatex report.tex} — вторая компиляция, включение библиографии
  \item \texttt{pdflatex report.tex} — третья компиляция, обновление ссылок
\end{enumerate}

Альтернативно можно использовать утилиту \texttt{latexmk}:

\begin{lstlisting}[language=bash, caption={Автоматическая компиляция с latexmk}]
latexmk -pdf -bibtex report.tex
\end{lstlisting}

\subsection{Результаты работы}

В результате выполнения лабораторной работы был создан документ с корректно оформленными библиографическими ссылками. Все цитирования в тексте автоматически связаны с соответствующими записями в списке литературы.

% ============================================
% АНАЛИЗ РЕЗУЛЬТАТОВ
% ============================================
\section{Анализ результатов}

\subsection{Преимущества использования biblatex}

По результатам выполнения работы можно выделить следующие преимущества использования biblatex:

\begin{enumerate}
  \item \textbf{Автоматизация} — список литературы формируется автоматически на основе цитирований в тексте
  \item \textbf{Единообразие} — все ссылки оформляются в едином стиле
  \item \textbf{Гибкость} — легко изменить стиль оформления, изменив один параметр
  \item \textbf{Переиспользование} — один файл \texttt{.bib} можно использовать в разных документах
  \item \textbf{Поддержка Unicode} — корректная работа с кириллицей и другими алфавитами
\end{enumerate}

\subsection{Сравнение с ручным оформлением}

В таблице \ref{tab:comparison} представлено сравнение автоматического и ручного оформления библиографии.

\begin{table}[htbp]
  \centering
  \caption{Сравнение методов оформления библиографии}
  \label{tab:comparison}
  \begin{tabular}{@{}lcc@{}}
    \toprule
    \textbf{Критерий} & \textbf{BibTeX/biblatex} & \textbf{Ручное оформление} \\
    \midrule
    Время на оформление & Низкое & Высокое \\
    Вероятность ошибок & Низкая & Высокая \\
    Единообразие стиля & Гарантировано & Требует внимания \\
    Изменение стиля & Одна строка кода & Переделка всего списка \\
    Переиспользование & Легко & Сложно \\
    Автоматическая сортировка & Да & Нет \\
    \bottomrule
  \end{tabular}
\end{table}

\subsection{Рекомендации по использованию}

На основе полученного опыта можно дать следующие рекомендации:

\begin{itemize}
  \item Использовать biblatex вместо устаревшего BibTeX для новых проектов
  \item Выбирать backend=biber для полной поддержки Unicode
  \item Хранить библиографию в отдельном файле \texttt{.bib} для переиспользования
  \item Использовать менеджеры библиографии (JabRef, Zotero) для управления большими базами
  \item Регулярно делать резервные копии файла \texttt{.bib}
  \item Проверять корректность записей перед финальной компиляцией
\end{itemize}

% ============================================
% ЗАКЛЮЧЕНИЕ
% ============================================
\section{Заключение}

В ходе выполнения лабораторной работы были изучены основные принципы работы с библиографией в системе LaTeX. Освоены инструменты BibTeX и biblatex, изучены различные типы библиографических записей и команды цитирования.

Практическая часть работы продемонстрировала эффективность использования автоматизированных средств управления библиографией. Создан полноценный документ с корректно оформленными ссылками на различные типы источников: книги, статьи, онлайн-ресурсы, материалы конференций.

Полученные навыки могут быть применены при подготовке курсовых работ, дипломных проектов, научных статей и других документов, требующих оформления библиографических ссылок по установленным стандартам.

Все поставленные задачи выполнены, цель работы достигнута.

% ============================================
% СПИСОК ЛИТЕРАТУРЫ
% ============================================
\newpage
\printbibliography[title={Список использованных источников}]

% Можно также разделить по типам:
% \printbibliography[type=book, title={Книги}]
% \printbibliography[type=article, title={Статьи}]
% \printbibliography[type=online, title={Интернет-ресурсы}]

% ============================================
% ПРИЛОЖЕНИЯ
% ============================================
\newpage
\appendix

\section{Пример файла библиографии}

Ниже приведен фрагмент файла \texttt{references.bib}:

\begin{lstlisting}[language=TeX, caption={Фрагмент references.bib}]
@book{knuth1984,
  author    = {Donald E. Knuth},
  title     = {The TeXbook},
  publisher = {Addison-Wesley Professional},
  year      = {1984},
  address   = {Reading, Massachusetts},
  isbn      = {0-201-13447-0},
  edition   = {1st}
}

@article{dijkstra1968,
  author  = {Edsger W. Dijkstra},
  title   = {Go To Statement Considered Harmful},
  journal = {Communications of the ACM},
  year    = {1968},
  volume  = {11},
  number  = {3},
  pages   = {147--148},
  doi     = {10.1145/362929.362947}
}

@online{latex-project,
  author  = {{The LaTeX Project}},
  title   = {LaTeX – A document preparation system},
  year    = {2024},
  url     = {https://www.latex-project.org/},
  urldate = {2024-01-15}
}
\end{lstlisting}

\section{Команды компиляции}

\subsection{Ручная компиляция}

\begin{lstlisting}[language=bash, caption={Последовательность команд компиляции}]
# Первая компиляция LaTeX
pdflatex report.tex

# Обработка библиографии
biber report

# Вторая компиляция LaTeX
pdflatex report.tex

# Третья компиляция для обновления ссылок
pdflatex report.tex
\end{lstlisting}

\subsection{Автоматическая компиляция}

\begin{lstlisting}[language=bash, caption={Использование latexmk}]
# Компиляция с автоматической обработкой библиографии
latexmk -pdf -bibtex report.tex

# Очистка вспомогательных файлов
latexmk -c

# Полная очистка (включая PDF)
latexmk -C
\end{lstlisting}

\section{Полезные ресурсы}

\subsection{Официальная документация}

\begin{itemize}
  \item The LaTeX Project: \url{https://www.latex-project.org/}
  \item CTAN — Comprehensive TeX Archive Network: \url{https://www.ctan.org/}
  \item Biblatex manual: \url{https://ctan.org/pkg/biblatex}
  \item Biber manual: \url{https://ctan.org/pkg/biber}
\end{itemize}

\subsection{Обучающие материалы}

\begin{itemize}
  \item Overleaf Learn LaTeX: \url{https://www.overleaf.com/learn}
  \item TeX Stack Exchange: \url{https://tex.stackexchange.com/}
  \item LaTeX.ru — русскоязычное сообщество: \url{https://www.latex.ru/}
\end{itemize}

\subsection{Инструменты}

\begin{itemize}
  \item JabRef — менеджер библиографии: \url{https://www.jabref.org/}
  \item Zotero — менеджер ссылок: \url{https://www.zotero.org/}
  \item Overleaf — онлайн редактор LaTeX: \url{https://www.overleaf.com/}
\end{itemize}

\end{document}

